\section{La fusión de razonamiento y acción: ReAct}

\subsection{El paradigma de Reasoning y el paradigma de Acting}

Antes de que apareciera ReAct, las aplicaciones de los LLM se dividían en dos paradigmas independientes:

\begin{itemize}
    \item \textbf{Reasoning} (como Chain-of-Thought): el LLM razona internamente paso a paso, lo que aumenta el test-time compute, pero no puede acceder a conocimiento externo ni a herramientas.
    \item \textbf{Acting} (como RAG, tool use): el LLM invoca el entorno externo (motores de búsqueda, calculadoras, API) para obtener conocimiento y capacidad de cómputo, pero carece de razonamiento para adaptarse a situaciones complejas.
\end{itemize}

Cada paradigma resuelve por separado un subproblema distinto de las preguntas y respuestas (problemas de razonamiento frente a problemas de conocimiento frente a problemas de cómputo), lo que produce métodos fragmentados (piecemeal).

\subsection{ReAct: un marco unificado de razonamiento y acción}

\begin{importantbox}{La idea central de ReAct}
Hacer que el LLM alterne entre generar un \textbf{pensamiento} (thought) y una \textbf{acción} (action), mientras el entorno devuelve una observación (observation). Basta con escribir un ejemplo de trayectoria como prompt que muestre cómo pensar y actuar para resolver la tarea, y el LLM puede imitar ese patrón.
\end{importantbox}

Flujo de trabajo concreto:
\begin{enumerate}
    \item El LLM genera un thought (por ejemplo, «necesito buscar la capitalización de mercado de estas tres empresas»)
    \item El LLM genera una action (por ejemplo, \texttt{Google{[}Apple market cap{]}})
    \item El entorno devuelve una observation (el resultado de la búsqueda)
    \item El thought, la action y la observation se agregan al contexto y el LLM continúa generando el siguiente paso
    \item Se repite hasta completar la tarea
\end{enumerate}

\begin{knowledgebox}{El beneficio mutuo entre razonamiento y acción}
\begin{itemize}
    \item \textbf{La acción ayuda al razonamiento}: obtiene información en tiempo real del exterior y ejecuta cálculos
    \item \textbf{El razonamiento ayuda a la acción}: planea estrategias según la situación actual y reajusta el plan ante anomalías. Por ejemplo, cuando la búsqueda devuelve el precio de la acción en lugar de la capitalización de mercado, el razonamiento puede indicar que el siguiente paso es buscar el número de acciones para calcularla de forma indirecta.
\end{itemize}
\end{knowledgebox}

\subsection{La generalidad de ReAct}

ReAct no se limita a las preguntas y respuestas: cualquier tarea que pueda convertirse en un «juego de texto» puede usar el marco de ReAct:

\begin{itemize}
    \item Navegación web, compras en línea
    \item Videojuegos (convirtiendo los píxeles en texto mediante image captioning)
    \item Simulación de tareas domésticas (ALFWorld)
    \item Ingeniería de software (SWE-bench)
\end{itemize}

\begin{warningbox}{Las limitaciones de un Acting Agent sin razonamiento}
En el juego ALFWorld, un agente puramente de acting (que no genera thought) repite la misma acción una y otra vez al encontrar una anomalía (por ejemplo, cuando el objeto buscado no está en la posición esperada), sin poder adaptarse. Al agregar la thinking action, el agente puede reflexionar sobre la situación actual y replanificar.
\end{warningbox}

\subsection{Resumen de la sección}

La aportación de ReAct no es solo un método concreto, sino la propuesta del paradigma del Reasoning Agent, que trata el razonamiento como una «acción interna» especial del agente. Este paradigma supera sistemáticamente tanto al razonamiento puro como a la acción pura.

